% Options for packages loaded elsewhere
\PassOptionsToPackage{unicode}{hyperref}
\PassOptionsToPackage{hyphens}{url}
%
\documentclass[
]{article}
\usepackage{amsmath,amssymb}
\usepackage{iftex}
\ifPDFTeX
  \usepackage[T1]{fontenc}
  \usepackage[utf8]{inputenc}
  \usepackage{textcomp} % provide euro and other symbols
\else % if luatex or xetex
  \usepackage{unicode-math} % this also loads fontspec
  \defaultfontfeatures{Scale=MatchLowercase}
  \defaultfontfeatures[\rmfamily]{Ligatures=TeX,Scale=1}
\fi
\usepackage{lmodern}
\ifPDFTeX\else
  % xetex/luatex font selection
\fi
% Use upquote if available, for straight quotes in verbatim environments
\IfFileExists{upquote.sty}{\usepackage{upquote}}{}
\IfFileExists{microtype.sty}{% use microtype if available
  \usepackage[]{microtype}
  \UseMicrotypeSet[protrusion]{basicmath} % disable protrusion for tt fonts
}{}
\makeatletter
\@ifundefined{KOMAClassName}{% if non-KOMA class
  \IfFileExists{parskip.sty}{%
    \usepackage{parskip}
  }{% else
    \setlength{\parindent}{0pt}
    \setlength{\parskip}{6pt plus 2pt minus 1pt}}
}{% if KOMA class
  \KOMAoptions{parskip=half}}
\makeatother
\usepackage{xcolor}
\usepackage[margin=2.5cm]{geometry}
\setlength{\emergencystretch}{3em} % prevent overfull lines
\providecommand{\tightlist}{%
  \setlength{\itemsep}{0pt}\setlength{\parskip}{0pt}}
\setcounter{secnumdepth}{-\maxdimen} % remove section numbering
\ifLuaTeX
\usepackage[bidi=basic]{babel}
\else
\usepackage[bidi=default]{babel}
\fi
\babelprovide[main,import]{australian}
% get rid of language-specific shorthands (see #6817):
\let\LanguageShortHands\languageshorthands
\def\languageshorthands#1{}
\usepackage{float}
\floatplacement{figure}{H}
\usepackage{newunicodechar}
\newunicodechar{α}{\ensuremath{\alpha}}
\newunicodechar{β}{\ensuremath{\beta}}
\newunicodechar{χ}{\ensuremath{\chi}}
\newunicodechar{δ}{\ensuremath{\delta}}
\newunicodechar{ε}{\ensuremath{\epsilon}}
\newunicodechar{η}{\ensuremath{\eta}}
\newunicodechar{μ}{\ensuremath{\mu}}
\newunicodechar{ρ}{\ensuremath{\rho}}
\newunicodechar{σ}{\ensuremath{\sigma}}
\newunicodechar{τ}{\ensuremath{\tau}}
\newunicodechar{Δ}{\ensuremath{\Delta}}
\newunicodechar{Σ}{\ensuremath{\Sigma}}
\ifLuaTeX
  \usepackage{selnolig}  % disable illegal ligatures
\fi
\IfFileExists{bookmark.sty}{\usepackage{bookmark}}{\usepackage{hyperref}}
\IfFileExists{xurl.sty}{\usepackage{xurl}}{} % add URL line breaks if available
\urlstyle{same}
\hypersetup{
  pdflang={en-AU},
  hidelinks,
  pdfcreator={LaTeX via pandoc}}

\author{}
\date{}

\begin{document}

\hypertarget{name-student-id}{%
\subsection{0. Name \& Student ID}\label{name-student-id}}

Michael Rideout - 36809268

\hypertarget{proposed-hypothesis}{%
\subsection{1. Proposed Hypothesis}\label{proposed-hypothesis}}

Among Australians aged 18--30 who do not own a home, those who view a
simulated social media feed of pessimistically framed housing posts will
report a lower perceived likelihood of owning a home within ten years
than those who view the same housing facts in neutral framing.

\hypertarget{null-hypothesis}{%
\subsection{2. Null Hypothesis}\label{null-hypothesis}}

Among Australians aged 18--30 who do not own a home, perceived
likelihood of owning a home within ten years will not differ between
those who view a simulated social media feed of pessimistically framed
housing posts and those who view the same housing facts in neutral
framing.

\hypertarget{independent-dependent-and-confounding-variables}{%
\subsection{3. Independent, Dependent and Confounding
Variables}\label{independent-dependent-and-confounding-variables}}

\begin{itemize}
\tightlist
\item
  \textbf{Independent:} post framing (pessimistic or neutral), randomly
  assigned.
\item
  \textbf{Dependent:} perceived likelihood (0--100\%) of owning a home
  within ten years.
\item
  \textbf{Confounding:} household income and savings; living arrangement
  (renting or with parents); capital-city versus regional location;
  prior weekly exposure to housing posts. All are measured before
  exposure; randomisation balances them.
\end{itemize}

\hypertarget{data-collection-100-words}{%
\subsection{4. Data Collection (\textless{} 100
words)}\label{data-collection-100-words}}

An online between-subjects experiment will recruit 128 valid completers
(64 per group). Voluntary-response convenience sampling is cheap and
reaches young non-homeowners online. Instagram and TikTok
advertisements, university notices and community groups link to a
Qualtrics survey screening age, residence and home ownership.
Participants answer confounder questions, are randomly allocated, view
eight mock posts with identical housing statistics across conditions,
and rate ownership likelihood, tone (manipulation check) and an
attention item before a debrief. Costs (AUD): incentives 640 (5 each),
advertising 250, piloting and contingency 110. Hours: stimuli/pilot 25,
ethics/setup 15, recruitment 25, analysis 15, reporting 20.

\hypertarget{statistical-tests-150-words}{%
\subsection{5. Statistical Tests (\textless{} 150
words)}\label{statistical-tests-150-words}}

Welch's independent-samples t-test will compare mean ownership
likelihood between groups. It suits a between-subjects design with two
independent groups and a continuous 0--100 outcome. Unlike Student's
t-test, it does not assume equal variances, which pessimistic framing
may widen by polarising responses. The test assumes independent
observations (one response per person, random allocation),
interval-level measurement and approximately normal scores within each
group; if Q--Q plots show floor or ceiling clustering, a Mann--Whitney U
test, which compares rank distributions, will replace it. The hypothesis
is directional because emotive, barrier-focused framing should make
identical facts seem less surmountable, but the test is two-sided (α =
.05) so an opposite effect is not missed. H0 is rejected if p
\textless{} .05; H1 is supported only if the pessimistic mean is lower.
The mean difference, 95\% CI and Cohen's d will be reported. Assuming a
medium effect (d = 0.5), 64 per group gives 80\% power.

\hypertarget{limitations-200-words}{%
\subsection{6. Limitations (\textless{} 200
words)}\label{limitations-200-words}}

A brief, researcher-made feed lacks the personalisation and repeated
exposure of real platforms, so any effect may not transfer to everyday
use; a multi-day field study with a browser-extension feed could test
this. Self-selected online volunteers probably over-represent digitally
engaged young people, which limits generalisation to all non-homeowners
aged 18--30; recruiting offline at TAFE campuses and workplaces would
reach less digitally engaged peers. An immediate single-item rating may
capture passing mood rather than a stable perception and invite demand
effects that overstate lasting change; a delayed follow-up with a
multi-item measure would show whether the effect persists.

\hypertarget{alternative-narrative-50-words}{%
\subsection{7. Alternative Narrative (\textless{} 50
words)}\label{alternative-narrative-50-words}}

``Housing pessimism reflects finances, not feeds'' could be manufactured
by recruiting only low-income capital-city renters, whose near-floor
ratings leave framing no room; measuring likelihood a week later, so
effects fade; and excluding pessimistic-condition participants who
passed the manipulation check. Equal near-floor group means would then
`show' framing is irrelevant.

\end{document}
